% 原创教学拓扑；左图不是实测脑连接，右图是有限前馈计算依赖。
\begin{tikzpicture}[x=1mm,y=1mm,font=\figurefont\footnotesize,
  cell/.style={circle,line width=1.1pt,draw=NNHidden,fill=NNHidden!18,minimum size=4mm,inner sep=0pt},
  othercell/.style={rectangle,line width=1.1pt,rounded corners=1.5mm,draw=NNHidden,fill=NNPanel,minimum size=4mm,inner sep=0pt},
  link/.style={nn legacy arrow,line width=.85pt}]
  \node[font=\figurefont\bfseries\small] at (23,52) {生物神经网络的局部示意};
  \node[font=\figurefont\bfseries\small] at (86,52) {前馈人工神经网络};
  \path[fill=NNInput!8,rounded corners=3mm] (1,26) rectangle (23,45);
  \path[fill=NNHidden!8,rounded corners=3mm] (27,15) rectangle (48,38);
  \path[fill=NNGradient!8,rounded corners=3mm] (3,1) rectangle (25,20);
  \foreach \i/\x/\y in {1/6/38,2/18/39,3/12/30,4/31/29,5/42/32,6/39/20,7/8/10,8/19/13}{
    \node[cell] (b\i) at (\x,\y) {};
  }
  \node[othercell] (b9) at (13,4) {};
  \foreach \i/\j in {1/2,2/3,3/1,4/5,5/6,6/4,7/8,8/9,9/7,3/4,6/8,7/3}{
    \draw[link] (b\i)--(b\j);
  }
  \draw[link] (b5) to[bend right=38] (b2);
  \foreach \i/\y in {1/8,2/23,3/38}{
    \node[cell,draw=NNInput,fill=NNInput!18] (x\i) at (65,\y) {};
  }
  \foreach \i/\y in {1/8,2/23,3/38}{\node[cell] (h\i) at (85,\y) {};}
  \foreach \i/\y in {1/15,2/31}{
    \node[cell,draw=NNOutput,fill=NNOutput!18] (o\i) at (105,\y) {};
  }
  \foreach \i in {1,2,3}{\foreach \j in {1,2,3}{\draw[link] (x\i)--(h\j);}}
  \foreach \i in {1,2,3}{\foreach \j in {1,2}{\draw[link] (h\i)--(o\j);}}
  \node[text=FigureMuted] at (65,0) {输入};
  \node[text=FigureMuted] at (85,0) {表示};
  \node[text=FigureMuted] at (105,0) {输出};
\end{tikzpicture}
